\section{Aspectos legais relevantes para a formalização}
\label{sec:legal}
\subsection{Registro na Junta Comercial e inscrição no CNPJ}
A sequência proposta começa com consulta de viabilidade do nome e do endereço, definição do objeto social e das atividades econômicas, elaboração do contrato social e protocolo na \textbf{JUCEMG}. O cadastro no CNPJ deve ser realizado conforme o fluxo integrado da \textbf{Redesim}, com eventuais inscrições municipais ou estaduais exigidas pela atividade (REDESIM, s.d.; JUCEMG, s.d.).

A consulta do nome empresarial na Junta não substitui a pesquisa nem o registro de marca no INPI. O contrato deverá identificar sócios, sede, objeto, capital, quotas, administração e regras essenciais de deliberação. A atividade real deverá corresponder ao objeto declarado; a empresa não deve ser registrada como clínica se apenas intermedeia serviços, nem ocultar prestação veterinária que de fato execute.

Os CNAEs deverão ser definidos com o contador a partir das receitas efetivas: software, licenciamento, assinaturas e eventual intermediação. Não se escolhe um código apenas para obter menor imposto. CNAE e objeto social influenciam licenciamento, tributação e obrigações profissionais.

\subsection{Licenças, alvarás e limites da atuação veterinária}
Mesmo uma operação digital deve verificar viabilidade do endereço, licenciamento municipal e classificação de risco. A eventual dispensa de alvará para baixo risco depende das regras aplicáveis e não dispensa cumprimento de normas sanitárias, urbanísticas ou de segurança (REDESIM, s.d.; JUCEMG, s.d.). A necessidade de inscrição estadual depende da existência de operações sujeitas ao ICMS; não se presume sua exigência para todo aplicativo.

As clínicas parceiras devem manter seus próprios registros profissionais e licenças exigíveis, incluindo responsabilidade técnica e requisitos sanitários conforme a atividade. O \projeto{} deve conferir documentação antes do credenciamento e estabelecer obrigação de atualização contratual.

Caso a plataforma passe a prestar serviços veterinários diretamente, deverá avaliar registro no CRMV, responsabilidade técnica e demais exigências profissionais. A regra da Lei nº 6.839/1980 considera a atividade básica ou o serviço prestado a terceiros; ser uma empresa de tecnologia não é, por si só, garantia de dispensa (Lei nº 6.839/1980).

\textbf{Triagem e orientação.} A interface deve separar conteúdo educativo de ato profissional e de atendimento clínico. A Resolução CFMV nº 1.465/2022 disciplina a telemedicina veterinária. Teleorientação e teletriagem não autorizam diagnóstico, solicitação de exames ou prescrição; atos privativos devem ser conduzidos por médico-veterinário dentro dos requisitos aplicáveis. Um questionário automatizado não deve diagnosticar, prescrever nem induzir o tutor a postergar atendimento urgente. A curadoria interna, por si só, não comprova autorização para prática clínica remota (CFMV, 2022).

\subsection{Regime tributário inicial}
Propõe-se \textbf{avaliar o Simples Nacional como primeira alternativa}, e não escolhê-lo de forma definitiva antes de classificar as atividades. O porte projetado, a estrutura de pessoa jurídica e as vedações da LC nº 123/2006 precisam ser examinados; o limite geral de receita do regime é R\$~4,8 milhões anuais. A opção não é automática após a inscrição no CNPJ (LC nº 123/2006; GOVERNO FEDERAL, s.d.).

{\small
\begin{longtable}{@{}L{3cm} L{6cm} L{5.8cm}@{}}
\caption{Regimes tributários a avaliar (Lei nº 9.718/1998; LC nº 123/2006).}\\
\toprule
\cabfundo \cab{Regime} & \cab{Características} & \cab{Aplicação proposta}\\
\midrule
\endfirsthead
\toprule
\cabfundo \cab{Regime} & \cab{Características} & \cab{Aplicação proposta}\\
\midrule
\endhead
\bottomrule
\endfoot
Simples Nacional & Recolhimento unificado de tributos abrangidos pelo regime; anexos e alíquotas dependem da atividade e receita. Há impedimentos societários e operacionais. & Primeira hipótese para o início, se elegível e economicamente favorável. Não equivale a imposto fixo de 6\% sobre toda receita.\\
Lucro Presumido & Bases presumidas de IRPJ e CSLL conforme atividade; demais tributos possuem regras próprias. Em regra, opção para receita no ano anterior até R\$~78 milhões, se não houver obrigação de Lucro Real. & Comparar em simulação com receitas segregadas, folha e margens. Não é necessariamente mais caro nem automaticamente adequado.\\
Lucro Real & IRPJ e CSLL sobre lucro fiscal ajustado; maior exigência de controles. Obrigatório em hipóteses legais e possível por opção. & Avaliar se houver impedimento aos demais regimes, crescimento ou vantagem demonstrada por simulação. Prejuízo contábil não elimina todos os tributos.\\
\end{longtable}}

Para serviços sujeitos ao \emph{fator R}, a relação entre folha e receita dos doze meses, nos termos legais, pode determinar tributação pelo Anexo III ou V. O patamar é 28\%; as alíquotas nominais iniciais são 6\% e 15,5\%, respectivamente. A aplicação depende da atividade, da faixa e da base correta; não basta dizer que ``sem pró-labore toda receita vai para 15,5\%''. A alíquota efetiva considera receita acumulada e parcela a deduzir (LC nº 123/2006).

A premissa tributária do modelo deve ser recalculada após a definição dos CNAEs e contratos. Valores recebidos de tutores, repasses às clínicas, comissões, assinaturas e patrocínios precisam ser contabilizados conforme sua natureza: não se pode deduzir repasses da receita tributável sem fundamento jurídico e documental. O contador deverá ainda conferir os efeitos e obrigações da transição da reforma tributária, especialmente para planejamento de 2027; não se acrescenta uma alíquota genérica de IBS/CBS ao orçamento do projeto (LC nº 214/2025).

\subsection{Contrato social e acordo entre sócios}
Além das cláusulas obrigatórias do contrato social, recomenda-se acordo de sócios compatível com ele, contemplando:
\begin{itemize}
\item \textbf{Aportes e administração:} quotas, prazos de integralização, poderes de assinatura, limites de gastos e aprovação de operações com partes relacionadas.
\item \textbf{Deliberações:} matérias reservadas e quóruns legais/contratuais. Quatro participações iguais podem gerar empate; prever negociação, mediação e solução final de impasses, sem contrariar quóruns obrigatórios.
\item \textbf{Entrada e saída:} preferência, condições para cessão de quotas, critérios de avaliação e apuração de haveres, prazo de pagamento, falecimento e exclusão nas hipóteses legais.
\item \textbf{Dedicação e remuneração:} funções, disponibilidade, pró-labore e distribuição de resultados, conforme caixa, escrituração e regras legais.
\item \textbf{Investimento e propriedade intelectual:} diluição, condições de conversão, sigilo, cessão de direitos e eventual vesting formalmente negociado.
\end{itemize}
Nenhuma dessas regras é apresentada como já acordada. O documento definitivo deve ser revisado por profissional habilitado (Código Civil, 2002; LC nº 182/2021).

\subsection{Marca, software e propriedade intelectual}
\begin{description}[style=nextline,leftmargin=0cm]
\item[Marca PetCuida.] Pesquisar anterioridades e requerer registro no INPI nas classes pertinentes. A exclusividade marcária decorre do registro validamente concedido, com vigência de dez anos a partir da concessão, renovável. Nome empresarial, domínio e perfil em redes sociais não substituem essa proteção (Lei nº 9.279/1996; INPI, s.d.-a).
\item[Software e conteúdo.] A Lei nº 9.609/1998 protege programas de computador independentemente de registro. O registro facultativo no INPI pode reforçar prova de autoria e titularidade. A lei prevê proteção por cinquenta anos, com marco inicial próprio. Formalizar os direitos de exploração econômica do código, design e materiais produzidos por fundadores e contratados; não presumir que todo trabalho acadêmico pertença automaticamente à futura empresa (Lei nº 9.609/1998; INPI, s.d.-b).
\item[Patentes.] Programa de computador em si, método comercial e ideia abstrata não são invenções patenteáveis pela Lei nº 9.279/1996. Eventual invenção implementada por computador requer análise técnica de requisitos; o aplicativo não deve ser apresentado como patenteável apenas por ser inovador (Lei nº 9.279/1996).
\item[Dependências e repositórios.] Inventariar bibliotecas, imagens, mapas e serviços de terceiros e cumprir licenças. Código publicado no GitHub não é automaticamente de domínio público; identificar titulares e condições de uso antes de incorporá-lo ao produto comercial.
\end{description}

\subsection{Proteção de dados pessoais e segurança}
A LGPD alcança dados dos tutores, profissionais e demais pessoas identificadas ou identificáveis. Informações do animal, isoladamente, não são dados pessoais humanos nem dados sensíveis de saúde do tutor; contudo, associadas ao cadastro, à localização e aos contatos, podem integrar tratamento de dados pessoais. Dados de saúde humana, se coletados, exigem análise específica (LGPD, 2018).

O \projeto{} deverá mapear finalidades e bases legais por operação, minimizar a coleta, definir retenção, publicar aviso de privacidade e oferecer canal para exercício de direitos. Consentimento não é base universal: execução contratual, obrigação legal e outras hipóteses precisam ser avaliadas separadamente. Compartilhamento com clínicas, ONGs e provedores deve ser documentado, incluindo funções de controlador/operador e transferências internacionais quando houver.

Devem ser implementados controle de acesso, autenticação, cópias de segurança, proteção dos registros e plano de resposta a incidentes, com avaliação de comunicação à ANPD e aos titulares nos termos aplicáveis. O regime simplificado para agentes de pequeno porte não isenta da LGPD e depende dos critérios da Resolução CD/ANPD nº 2/2022, incluindo restrições para tratamento de alto risco (LGPD, 2018; ANPD, 2022).

\subsection{Consumidores, parceiros e créditos solidários}
Termos de uso devem apresentar com clareza preços, cobrança recorrente, cancelamento, suporte, limites do serviço e identificação do fornecedor. Nas contratações pela internet, observar o CDC e o Decreto nº 7.962/2013, inclusive o direito de arrependimento quando aplicável. Cláusula de isenção não afasta responsabilidade legal nem corrige oferta enganosa (CDC, 1990; Decreto nº 7.962/2013).

O contrato com clínicas deverá disciplinar credenciamento, tabela social, qualidade, obrigações profissionais, agenda, repasses, cancelamentos, notas fiscais, atendimento de reclamações e tratamento de dados. O contrato com patrocinadores deve especificar destinação dos recursos, prestação de contas e autorização de uso de marcas.

Para os créditos solidários, formalizar emissão, utilização, validade, limites, cancelamento e liquidação, com transparência sobre eventual indisponibilidade e sem promessa de cobertura assistencial. O modelo prevê liquidação parcial em dinheiro e fundo de repasse; isso exige conciliação e contratos, mas não demonstra garantia financeira já constituída. Se o desenho evoluir para carteira, recursos de terceiros, crédito financeiro, seguro ou cobertura de riscos, deve haver avaliação regulatória específica antes do lançamento.

A Lei nº 9.608/1998 define serviço voluntário como atividade não remunerada prestada por pessoa física a entidade pública ou instituição privada sem fins lucrativos que atenda aos objetivos previstos na lei, mediante termo de adesão. Como a estrutura proposta para o PetCuida é uma sociedade empresária, não se deve presumir que esse regime se aplique diretamente à operação. Créditos resgatáveis por serviços podem representar contrapartida econômica e gerar questionamentos sobre a natureza remunerada da atividade e eventual vínculo de trabalho, conforme a forma concreta da prestação. Se houver modelo voluntário juridicamente cabível, recomenda-se termo de adesão, atuação eventual e sem subordinação, e créditos restritos a serviços de cuidado animal, sem conversão em salário. Essas medidas não eliminam o risco; a estrutura deve ser validada por assessoria jurídica antes do piloto (Lei nº 9.608/1998).

\subsection{Menores de idade e acesso ao aplicativo}
Como o aplicativo pode ser acessado por crianças e adolescentes, o tratamento de dados deverá observar seu melhor interesse, conforme o art.~14 da LGPD. Para serviços direcionados a esse público ou de acesso provável por ele, o Estatuto Digital da Criança e do Adolescente prevê proteção prioritária, elevado nível de privacidade e segurança, supervisão parental e aferição etária conforme os deveres aplicáveis a cada serviço (LGPD, 2018; Estatuto Digital da Criança e do Adolescente, 2025).

Como regra proposta, a assinatura será contratada por pessoa maior de 18 anos ou, se o usuário for menor, mediante representação ou assistência do responsável legal, conforme a idade e a natureza do ato. O mecanismo de verificação de idade, a autorização necessária e os controles parentais aplicáveis ao PetCuida deverão ser definidos com assessoria jurídica, observando a minimização de dados e as características efetivas do serviço (Código Civil, 2002; LGPD, 2018; Estatuto Digital da Criança e do Adolescente, 2025).